% !TEX root = ../ICRA19Hybrid.tex

\section{APPROACH}
\label{sec:approach}
% ===============================================================
Now we introduce an algorithm to efficiently solve the problem defined in section \ref{sub:problem_formulation}. The algorithm first solves for velocity commands, during which the dimensions and directions of both velocity control and force control are also determined. Then we fix the directions and solve for force controlled actions.

\subsection{Solve for Velocity Controlled Actions} % (fold)
\label{sub:solve_for_velocity_controlled_actions}
In this section, we design the velocity command (solve for $n_{af}, n_{av}, T$ and $w_{av}$), so as to satisfy all the velocity-level conditions.
We use a $n_{av}\times n$ selection matrix $S_{av}$ to select the velocity commands out of the generalized variables: $w_{av}=S_{av}w$. Equations of interest to this section are:
(Use $\dot q=\Omega v, w=Tv$, omitting argument $q$)
\begin{itemize}
  \item Holonomic natural constraint $J_\Phi\Omega v=0$. Denote $N=J_\Phi\Omega$, the constraint becomes $Nv=0$;
  \item Goal condition $Gv=b_G$;
  \item Velocity command $S_{av}Tv=w_{av}$. Denote $C=S_{av}T$, $b_C=w_{av}$, rewrite the velocity command as $Cv=b_C$.
\end{itemize}
Denote the solution set of each equation as $Sol(N), Sol(G)$ and $Sol(C)$. We need to design the velocity command $C,\ b_C$ such that the resulted solution space (the solution set of natural constraints and velocity commands) becomes \textit{a non-empty subset of the desired generalized velocities} (the solution set of natural constraints and goal condition):
\begin{equation}
Sol(N\&C)\ \in\ Sol(N\&G)
\end{equation}

\subsubsection{Determine dimensions of velocity control} % (fold)
\label{subsub:determine_dimensions_of_velocity_control}
Denote $r_N={\rm rank}(N), r_{NG}={\rm rank}(\left[ {\begin{array}{*{20}{c}}{N }\\G\end{array}}\right])$. The minimum number of independent velocity control we must enforce is
\begin{equation}
\label{eq:nav_min}
  {n^{\min}_{av}} = r_{NG}-r_N.
\end{equation}
This condition makes sure the dimension of $Sol(N\&C)$ is smaller or equal to the dimension of $Sol(N\&G)$, so that their containing relationship becomes possible. Physically it means the velocity commands need to reduce enough degree-of-freedoms from the null space of the natural constraints. The maximum number of independent velocity commands we can enforce is
\begin{equation}
\label{eq:nav_max}
  n^{\max}_{av} = n - r_N={\rm Dim}({\rm null}(N)),
\end{equation}
where ${\rm null}(N)$ denotes the null space of $N$. This condition ensures the system will not be overly constrained to have no solution.
We choose the minimal number of necessary velocity constraints:
\begin{equation}
\label{eq:choice_of_nav}
n_{av}=n^{\min}_{av} =r_{NG}-r_N.
\end{equation}
This choice makes it easier for the system to avoid jamming. As will be shown in the next section, it also leaves more space for solving force controlled actions.


\subsubsection{Solve for directions and magnitude} % (fold)
\label{subsub:solve_for_directions_and_magnitude}
With our choice of $n_{av}$, we know ${\rm rank}([N; C])={\rm rank}([N; G])$. Then the condition $Sol(N\&C)\ \in\ Sol(N\&G)$ implies
\begin{equation}
\label{eq:sol_NC_equals_sol_NG}
  Sol(N\&C)\ =\ Sol(N\&G),
\end{equation}
\ie the two linear systems share the same solution space. (\ref{eq:sol_NC_equals_sol_NG}) can be achieved by firstly choosing $C$ such that the homogeneous linear systems $\left[ {\begin{array}{*{20}{c}}N\\C\end{array}} \right]v = 0$ and $\left[ {\begin{array}{*{20}{c}}N\\G\end{array}} \right]v = 0$
become equivalent (share the same solution space). Compute a basis for the solution of $[N^T G^T]^Tv=0$: $[\sigma_1,...,\sigma_{n-r_{NG} }]$, then we just need to ensure $C$ satisfies:
\begin{equation}
\label{eq:constraints_on_C}
  C\sigma_i=0,\ \  i=1,...,n-r_{NG}
\end{equation}
Then we can compute $b_C$ from any specific solution of $\{Nv=0, Gv=b_G\}$. The original non-homogeneous systems then become equivalent.

Beside equation (\ref{eq:constraints_on_C}), we have a few more requirements/preferences on $C$ based on the discussions in section \ref{sub:velocity_controlled_actions_and_holonomic_constraints}:
\begin{itemize}
  \item Rows of $C$ must be linearly independent from each other. And we prefer to have them as orthogonal to each other as possible.
  \item Each row of $C$ is also independent from rows of the holonomic natural constraint $N$. We prefer to pick the rows as close to ${\rm null}(N)$ as possible.
\end{itemize}
To solve for $C$, denote $c^T\in \mathbb{R}^{1\times n}$ as any row in $C$. From $C=S_{av}T$ we know the first $n_u$ columns in $C$ are zeros, rewrite this and equation (\ref{eq:constraints_on_C}) as a linear constraint on $c$:
\begin{equation}
\label{eq:constraint_on_c}
  \left[ {\begin{array}{*{20}{c}}{\begin{array}{*{20}{c}}{\sigma _1^T}\\ \vdots \\{\sigma _{{n} - {n_{NG}}}^T}\end{array}}\\{\left[ {\begin{array}{*{20}{c}}{{I_{{n_u}}}}&{{{\bf{0}}_{{n_u} \times {n_a}}}}
  \end{array}} \right]}\end{array}} \right]c  = \left[ {\begin{array}{*{20}{c}}
  0\\ \vdots \\0\end{array}} \right]
\end{equation}
Its solution space has dimension of $n_c = n_a - n + r_{NG} = r_{NG} - n_u$. Since we need $n_{av}$ independent constraints, we require $n_c=r_{NG}-n_u \ge n_{av}=r_{NG}-r_N$, which gives $r_N\ge n_u $, \ie
\begin{equation}
\label{eq:dimensional_requirement}
r_N+n_a\ge n.
\end{equation}
For our method to work, (\ref{eq:dimensional_requirement}) says it must be possible for the actions and constraints to fully constrain the system.
Denote matrix $\mathbb{B}_c=[c^{(1)}\ \cdots\ c^{(n_c)}]$ as a basis of the solution space of equation (\ref{eq:constraint_on_c}). Denote ${\rm Null}(N)$ as a basis of ${\rm null}(N)$.
We can find a $C$ that satisfies all the conditions by solving the following optimization problem:
\begin{equation}
\label{eq:optimization_for_velocity}
\begin{array}{l}
\mathop {\min }\limits_{{{\bf{k}}_1}, \cdots ,{{\bf{k}}_{{n_{av}}}}} \sum\limits_{i \ne j} {||{c_i^T}c _j|{|}}  - \sum\limits_i {||{\rm Null}(N)^T c _i|{|}} \\
s.t.\;\;\,\;\;\,{c_i^T}c_i = 1,\;\;\,\forall i\\
\;\;\,\;\;\,\;\;\,\;\;\,{c _i} = \mathbb{B}_c{{\bf{k}}_i},\;\;\,\forall i
\end{array}
\end{equation}
The best velocity constraint $C^*=(\mathbb{B}_c\;[{\bf k_1\;\;\;...\;\; k_{n_{av}}}])^T$. The optimization problem (\ref{eq:optimization_for_velocity}) is non-convex because of the unit length constraint $c_i^Tc_i=1$. However, we can solve the problem numerically by projecting the solution back to the constraint after each gradient update:
\begin{enumerate}
  \item Start from a random ${\bf k=[{\bf k_1\;\;\;...\;\; k_{n_{av}}}]}$;
  \item Perform a gradient descent step: ${\bf k}\gets {\bf k} - t\nabla f$;
  \item Projection: ${\bf k_i}\gets \frac{\bf k_i}{||\mathbb{B}_c{\bf k_i}||},\;\;\;\;\forall i$;
  \item Repeat from step two until convergence.
\end{enumerate}
Here $\nabla f$ is the gradient of the cost function reshaped to the same size as $\bf k$. $t=10$ is a step length. In practice, we run the projected gradient descent algorithm above with $N_s$ different initializations to avoid bad local minima.

After obtaining $C^*$, we know the last $n_{av}$ rows of $R_a$. Denote the last $n_a$ columns of $C^*$ as $R_{C^*}$, we can expand it into a full rank $R_a$:
\begin{equation}
\label{eq:expand_R_a}
{R_a} = \left[ {\begin{array}{*{20}{c}}
{{\rm Null}{{({R_{C^*}})}^T}}\\
{{R_{C^*}}}
\end{array}} \right],
\end{equation}
it encodes the axes of the force controlled directions. Then we have $T=\rm{diag}(I_u, R_a)$.
The procedures are summarized in algorithm \ref{alg:solve_for_velocity}.
\begin{algorithm}[h]
\caption{Solve for velocity controlled actions}  \label{alg:solve_for_velocity}
\begin{algorithmic}[1]
% \renewcommand{\algorithmicrequire} {\textbf{Input :} }
% \REQUIRE $\qobj\ui,\qobj\uf,\qgrp\uf,\P$, number of samples $N_s$.
\STATE Check condition (\ref{eq:dimensional_requirement}) for feasibility.
\STATE Compute $n_{av}$ from equation (\ref{eq:choice_of_nav}).
\STATE Compute a basis of $[N^T G^T]^Tv=0$, plug in equation (\ref{eq:constraint_on_c}) and compute a basis $\mathbb{B}_c$.
\STATE Sample $N_s$ sets of coefficients ${\bf k}\in\mathbb{R}^{n_c\times n_{av}}$
\FOR {each sample $\bf k$}
\STATE Solve the optimization problem (\ref{eq:optimization_for_velocity}).
\STATE Compute $C=(\mathbb{B}_c{\bf k})^T$ from the solution.
\STATE Compute the cost of $C$ from equation (\ref{eq:optimization_for_velocity}).
\ENDFOR
\STATE Pick the $C^*$ with lowest cost.
\STATE Use equation (\ref{eq:expand_R_a}) to compute $R_a$. Then $T=\rm{diag}(I_u, R_a)$.
\STATE Compute one solution $v^*$ for $Nv=0, Gv=b_G$.
\STATE Compute $w_{av}=b_C=C^*v^*$.
\end{algorithmic}
\end{algorithm}

% subsection solve_for_velocity_controlled_actions (end)
\subsection{Solve for Force Controlled Actions} % (fold)
\label{sub:solve_for_force_controlled_actions}
Next we compute the force command (solve for $\eta_{af}$) so as to satisfy all the force-level requirements.
Equations of interest to this section:
(Use $\eta=Tf$, omitting argument $q$)
\begin{itemize}
  \item Newton's second law: express (\ref{eq:newton_second_law}) in the transformed generalized force space:
    \begin{equation}
      \label{eq:newton_second_law_lambda_eta}
      T\Omega^T(q)J^T_\Phi(q)\lambda + \eta + TF = 0.
    \end{equation}
  \item Guard conditions: express them as constraints on $\lambda, \eta$:
    \begin{equation}
      \label{eq:guard_condition_lambda_eta_inequality}
      \Lambda\left[ {\begin{array}{*{20}{c}}\lambda \\f\end{array}} \right] = [{\Lambda_\lambda }\;\;\,{\Lambda_f}{T^{ - 1}}]\left[ {\begin{array}{*{20}{c}}\lambda \\\eta \end{array}} \right] \le {b_\Lambda}.
    \end{equation}
    \begin{equation}
      \label{eq:guard_condition_lambda_eta_equality}
      \Gamma\left[ {\begin{array}{*{20}{c}}\lambda \\f\end{array}} \right] = [{\Gamma_\lambda }\;\;\,{\Gamma_f}{T^{ - 1}}]\left[ {\begin{array}{*{20}{c}}\lambda \\\eta \end{array}} \right] = {b_\Gamma}.
    \end{equation}
\end{itemize}
The unknowns are the force variables $\lambda, \eta$. Remember $\eta=[\eta_u, \eta_{af}, \eta_{av}]$. All the equality constraints ((\ref{eq:newton_second_law_lambda_eta}), (\ref{eq:guard_condition_lambda_eta_equality}) and our choice of $\eta_{af}$) will determine the value of all the forces; we need to make sure the resulted forces satisfy the inequality constraints (\ref{eq:guard_condition_lambda_eta_inequality}).
%
Remember also $f_u=0$. Express it as $Hf=HT^{-1}\eta=0$. Combine $HT^{-1}\eta=0$, (\ref{eq:newton_second_law_lambda_eta}) and (\ref{eq:guard_condition_lambda_eta_equality}) into one constraint:
\begin{equation}
  \label{eq:newton_all}
  \left[ {\begin{array}{*{20}{c}}{{0}}&{{H}{T^{ - 1}}}\\{T{\Omega ^T}J_\Phi ^T}&I\\\Gamma_\lambda&\Gamma_fT^{-1}\end{array}} \right]\left[ {\begin{array}{*{20}{c}}\lambda \\\eta \end{array}} \right] = \left[ {\begin{array}{*{20}{c}}{\bf{0}}\\{ - TF}\\b_\Gamma\end{array}} \right].
\end{equation}
Due to the limitation of rigid body modeling, the free forces may not have a unique solution given a force action $\eta_{af}$. Denote the free forces as $f_{free} = [\lambda^T\; \eta_u^T\; \eta_{av}^T]^T$, rewrite the constraints (\ref{eq:newton_all}) and move $\eta_{af}$ to the right hand side, we find one solution for $f_{free}$ by penalizing the sum-of-squares norm of the free forces:
\begin{equation}
  \label{eq:QP_newton_law}
  \begin{array}{l}\mathop {\min }\limits_{f_{free}} \;\;\,f_{free}^Tf_{free}\\s.t.\;\;\,\;\;\,{M_{free}}f_{free} = \left[ {\begin{array}{*{20}{c}}
  {\bf{0}}\\{ - TF}\\b_\Gamma\end{array}} \right] - {M_{{\eta _f}}}{\eta _f}.\end{array}
\end{equation}
This is a quadratic programming (QP) problem. Denote $f^*_{free}$ as the dual variables of $f_{free}$, the KKT condition says the solution to the QP can be found by solving the following linear system:
\begin{equation}
  \label{eq:KKT}
  \left[ {\begin{array}{*{20}{c}}{2I}&{M_{free}^T}\\{{M_{free}}}&{\bf 0}\end{array}} \right]\left[ {\begin{array}{*{20}{c}}f_{free}\\f_{free}^*\end{array}} \right] = \left[ {\begin{array}{*{20}{c}}{\bf 0}\\{\left[ {\begin{array}{*{20}{c}}{\bf{0}}\\{ - TF}\end{array}} \right] - {M_{{\eta _f}}}{\eta _f}}\end{array}} \right].
\end{equation}
This linear system uniquely determines the free forces given force action $\eta_{af}$. Rewrite it as
\begin{equation}
  \label{eq:KKT_final}
  \left[ {\begin{array}{*{20}{c}}{2I}&{M_{free}^T}&{\bf{0}}\\{{M_{free}}}&{\bf{0}}&{{M_{{\eta _f}}}}\end{array}} \right]\left[ {\begin{array}{*{20}{c}}f_{free}\\f_{free}^*\\\eta_{af}\end{array}} \right] = \left[ {\begin{array}{*{20}{c}}{\bf{0}}\\{ - TF}\end{array}} \right].
\end{equation}
This linear equation encodes the unique solution for Newton's law. Finally we solve (\ref{eq:KKT_final}) together with guard conditions (\ref{eq:guard_condition_lambda_eta_inequality}) to compute all forces. The procedure is summarized in algorithm \ref{alg:solve_for_force}.
\begin{algorithm}[h]
\caption{Solve for force controlled actions}  \label{alg:solve_for_force}
\begin{algorithmic}[1]
\STATE From Newton's laws, write down $M_{free}, M_{\eta_{af}}$ in (\ref{eq:QP_newton_law}).
\STATE Write down coefficient matrices for equation (\ref{eq:KKT_final}).
\STATE Solve the linear programming problem (\ref{eq:guard_condition_lambda_eta_inequality})(\ref{eq:KKT_final}) for $\eta_{af}$.
\end{algorithmic}
\end{algorithm}

% subsection solve_for_force_controlled_actions (end)